\shrink[1]
\section{Evaluation \label{sec:evaluation}}

\inred{To be written by the team.}

\shrink[1]
\subsection{Overview \label{sec:evaluation:overview}}

% No More LLMs

We evaluate \ToolName{} along the following axes. First, we present our main evaluation. More precisely, we evaluate whether the joint search for rewriting rules and rotation key selection is better than separate search (Sec.~\ref{sec:evaluation:joint-vs-separate}), we also evaluate whether our proposed approach protects the RL agent from exceeding the noise budget (Sec.~\ref{sec:evaluation:constrained-vs-unconstrained}). Second, we compare \ToolName{} to state-of-the-art compilers including CHEHAB RL~\cite{ChehabRL} and Coyote~\cite{Malik2023Coyote} (Sec.~\ref{sec:evaluation:comparison-with-sota}). Third, we provide an ablation study of the main components of our proposed approach (Sec.~\ref{sec:evaluation:ablations}). Additional experiments are discussed briefly while their details are deferred to the Supplementary Material.



\shrink[1]
\subsection{Experimental Setup \label{sec:evaluation:setup}}

% No More LLMs

\shrink[1]
\subsubsection{Benchmarks \label{sec:evaluation:setup:benchmarks}}

We evaluate \ToolName{} on three benchmark suites. These are the same benchmark suites used in the CHEHAB RL paper~\cite{ChehabRL}. For each benchmark in the benchmark suites, we vary the input size to evaluate the scalability of the compiler. The larger the input is, the larger the program is, since FHE code is fully unrolled. The three benchmark suites are:

% \shrink[0.5]
\paragraph{Porcupine Benchmark Suite} This is a set of kernels commonly used in linear algebra, machine learning, and image processing. They were used to evaluate Porcupine~\cite{Cowan2021Porcupine}. They include image processing filters: Box Blur, the x and y gradients of an image (\emph{Gx} and \emph{Gy}), and Robert Cross (\emph{Rob. Cross.}). The Porcupine benchmark suite also includes Dot Product, Hamming Distance (\emph{Hamm. Dist.}), L2 Distance, Linear Regression (\emph{Linear Reg.}), and Polynomial Regression (\emph{Poly. Reg.}), which are building blocks for ML applications. We include all of the benchmarks that were used to evaluate Porcupine.

% \shrink[0.5]
\paragraph{Coyote Benchmark Suite} It includes kernels that were used to evaluate Coyote~\cite{Malik2023Coyote}. These kernels are: 1) matrix multiplication (\emph{Mat. Mul.}); 2) a sorting algorithm (\emph{Sort}) that sorts a list of elements; This sorting algorithm is an unstructured code that implements sorting using a tree (as described by Malik et al.~\cite{Malik2023Coyote}).
3) \emph{Max} kernel. This is an unstructured code that finds the maximum element in a list.
Coyote also includes three other kernels, which are the dot product, the L2 distance, and a convolution, but these are common with the Porcupine benchmark suite, and therefore, we do not evaluate them again.
All of the benchmarks that were used to evaluate Coyote are included in our benchmark suite.

\shrink[0.25]
\paragraph{Randomly Generated Irregular Polynomials} The Coyote compiler also uses a set of randomly generated unstructured polynomials for evaluation. We use the same methodology for evaluation. This set of polynomials is a stress test to further investigate the ability of the compiler to vectorize code in the absence of a regular structure in the code. Several polynomials are generated as random arithmetic expression trees. In these polynomials, a polynomial is named tree-X-Y-Z, where high values for X-Y indicate that the polynomial tree is full, complete, and well-balanced, making it easier to vectorize, while low values for X-Y indicate that the tree is sparse and imbalanced, making it harder to vectorize. Z indicates the depth of the tree of the polynomial.\ifshowextra We provide more details about these polynomials in Appendix~\ref{irregularPolynomials}.\fi



\paragraph{Baselines}

We compare code generated by \ToolName{} with code generated by CHEHAB RL~\cite{ChehabRL} and code produced by the Coyote~\cite{Malik2023Coyote} compiler.
We limit our evaluation to FHE compilers that support unstructured code.
Porcupine~\cite{Cowan2021Porcupine} is excluded as it is not publicly available (we reached out to its authors but could not obtain a copy).


\paragraph{Evaluation Metrics}

We report two families of metrics: metrics related to the \emph{quality of programs generated} and metrics related to \emph{the compilation time}.

To evaluate the quality of programs generated by \ToolName{}, we report the execution time of the generated circuit, the number of homomorphic operations, the depth, the multiplicative depth, and noise (noise accumulated by the benchmark, which we want to minimize).\ifshowextra The definition of the noise metric and details about how we measure it are in Appendix~\ref{noise}.\fi To evaluate the compilation time, we report the end-to-end compilation time for \ToolName{} as well as the baselines.


\shrink[1]
\subsubsection{Environment Setup \label{sec:evaluation:setup:env}}

% Dalma.
\paragraph{Hardware}
We use the following machine for evaluation and training: a dual-socket Intel Xeon E5-2680 v4 with 28 cores (14 cores per socket, 2.40 GHz) and 105 GB of RAM running CentOS Linux 8.



\paragraph{\ToolName{} and Baseline Setup}

We follow the same setup used in CHEHAB RL~\cite{ChehabRL}. The encryption parameters for all the compilers are set as follows: n = 16384 (polynomial modulus degree) and q is the default coefficient modulus proposed by the SEAL library for a standard 128-bit security level.
We disable the CHEHAB RL pass for automatic rotation key selection~\cite{ChehabRL} (this is disabled for CHEHAB RL and for our system since it uses the CHEHAB compilation pipeline, since Coyote does not have a similar pass.
We also disable the use of the blocking technique \cite{Malik2023Coyote} for all the compilers, since they all perform blocking in the same way. Blocking is a technique that involves vectorizing smaller kernels separately and then composing the resulting vectorized programs. While this technique enables the two compilers to scale to larger programs, our goal is to compare the fundamental algorithms proposed by each compiler on a single block. Since the blocking method is identical in all the compilers, comparing on a single block will better illustrate the fundamental differences between the two approaches. 
We also use the same SEAL library in all the compilers (SEAL version 4.1). We run all the experiments 30 times and report the median execution times and compilation times. The compilation timeout for all compilers was 7200 seconds.


\subsection{RL Agent Training} \label{sec:evaluation:agent_training}
We train our RL agent for 5 million timesteps (96.5 hours). More details about \ifshowextra the training and \fi the hyperparameters in
\ifshowextra Appendix~\ref{training} and \fi Table~\ref{tab:ppo_hyperparameters}.

\input{tables/hyperparameters.tex}

\shrink[1]

\subsection{Main Results \label{sec:evaluation:main-results}}

% Replaces the empty \subsubsection{Unified CMO-RL vs. Original CHEHAB RL} in 50-evaluation.tex.
% Needs in util/packages.tex:  \usepackage{pgfplots} \usepgfplotslibrary{groupplots} \pgfplotsset{compat=1.17}
% All numbers in \inred{} are PLACEHOLDERS, not results.

\subsubsection{Unified CMO-RL vs. Original CHEHAB RL
\label{sec:evaluation:unified-vs-original}}

We compare \ToolName{} with the original CHEHAB RL~\cite{ChehabRL}. CHEHAB RL optimizes execution time only and ignores both rotation keys and noise. To give it a way to reduce key memory, we enable its rotation key selection pass for this experiment only. This pass runs after the RL agent and decomposes rotations into power-of-two rotations to reduce the number of keys (Sec.~\ref{sec:evaluation:setup:env}). We fix the noise budget to $B = \inred{230}$ bits, a budget that CHEHAB RL exceeds on most benchmarks. We fix the preference to $\vec{w} = [0.9, 0.1]^{\top}$, which co-optimizes speed and key memory. We report the raw output of the agent, without safety rollback (Sec.~\ref{sec:safety_rollback}).

\begin{figure}[t]
  \centering
  \pgfplotsset{exp1/.style={
      ybar, bar width=3.5pt, width=\linewidth, height=3.0cm, enlarge x limits=0.04,
      symbolic x coords={DP16,DP32,L2-16,L2-32,HD16,HD32,PR8,LR8,BB4,MM4,Gx4,Max4,Sort3,T100-50-5},
      xtick=data, xticklabel style={font=\tiny, rotate=30, anchor=east},
      yticklabel style={font=\tiny}, ylabel style={font=\scriptsize, align=center},
      ymajorgrids, grid style={gray!25},
      legend style={font=\scriptsize, draw=none, legend columns=-1, at={(0.5,1.02)}, anchor=south}}}
  \begin{tikzpicture}
  \begin{groupplot}[group style={group size=1 by 3, vertical sep=0.3cm, xticklabels at=edge bottom}]
  \nextgroupplot[exp1, xticklabel=\empty, ylabel={Noise / $B$}, ymin=0.7, ymax=1.25,
                 legend entries={CHEHAB RL + key sel., \ToolName{}}]
    \addplot[fill=gray!55, draw=gray!80] coordinates {(DP16,1.048) (DP32,1.122) (L2-16,1.026) (L2-32,1.139) (HD16,1.074) (HD32,1.178) (PR8,0.952) (LR8,0.817) (BB4,0.891) (MM4,1.096) (Gx4,0.930) (Max4,1.013) (Sort3,1.039) (T100-50-5,1.157)};
    \addplot[fill=blue!60, draw=blue!80] coordinates {(DP16,0.974) (DP32,0.987) (L2-16,0.961) (L2-32,0.996) (HD16,0.983) (HD32,0.991) (PR8,0.922) (LR8,0.817) (BB4,0.883) (MM4,0.983) (Gx4,0.913) (Max4,0.978) (Sort3,0.996) (T100-50-5,0.991)};
    \draw[red, dashed, thick] (rel axis cs:0,0.5454) -- (rel axis cs:1,0.5454);
  \nextgroupplot[exp1, xticklabel=\empty, ylabel={Speedup}, ymin=0.9, ymax=1.1, bar width=5pt]
    \addplot[fill=blue!60, draw=blue!80] coordinates {(DP16,1.017) (DP32,1.022) (L2-16,0.998) (L2-32,1.023) (HD16,1.030) (HD32,1.063) (PR8,1.019) (LR8,1.000) (BB4,1.043) (MM4,1.028) (Gx4,1.017) (Max4,1.019) (Sort3,1.010) (T100-50-5,1.013)};
    \draw[black] (rel axis cs:0,0.5) -- (rel axis cs:1,0.5);
  \nextgroupplot[exp1, ylabel={Keys (MB)}, ymin=0, ymax=140]
    \addplot[fill=gray!55, draw=gray!80] coordinates {(DP16,90) (DP32,108) (L2-16,90) (L2-32,108) (HD16,90) (HD32,126) (PR8,36) (LR8,36) (BB4,54) (MM4,72) (Gx4,54) (Max4,90) (Sort3,90) (T100-50-5,108)};
    \addplot[fill=blue!60, draw=blue!80] coordinates {(DP16,54) (DP32,72) (L2-16,54) (L2-32,72) (HD16,72) (HD32,72) (PR8,36) (LR8,18) (BB4,36) (MM4,54) (Gx4,36) (Max4,72) (Sort3,72) (T100-50-5,72)};
  \end{groupplot}
  \end{tikzpicture}
  \caption{\ToolName{} vs.\ CHEHAB RL with rotation key selection ($B=\inred{230}$ bits, $\vec{w}=[0.9,0.1]^{\top}$). Top: noise normalized by $B$ (bars above the dashed line violate the budget). Middle: speedup of \ToolName{} (higher is better). Bottom: rotation key size (lower is better). \inred{Placeholder data.}}
  \label{fig:exp-unified-vs-original}
\end{figure}

\begin{table}[t]
  \centering
  \caption{Noise (bits), budget violation, execution time (ms), and number of rotation keys ($B=\inred{230}$ bits, $\vec{w}=[0.9,0.1]^{\top}$). \inred{Placeholder data.}}
  \label{tab:exp-unified-vs-original}
  \fontsizefortables
  \setlength{\tabcolsep}{3.5pt}
  \begin{tabular}{@{}l rcrr rcrr@{}}
    \toprule
    & \multicolumn{4}{c}{CHEHAB RL + key sel.} & \multicolumn{4}{c}{\ToolName{}} \\
    \cmidrule(lr){2-5}\cmidrule(l){6-9}
    Benchmark & Noise & OK & Time & Keys & Noise & OK & Time & Keys \\
    \midrule
    Dot Product 16  & \inred{241} & \no  & \inred{412}  & \inred{5} & \inred{224} & \yes & \inred{405}  & \inred{3} \\
    Dot Product 32  & \inred{258} & \no  & \inred{655}  & \inred{6} & \inred{227} & \yes & \inred{641}  & \inred{4} \\
    L2 Distance 16  & \inred{236} & \no  & \inred{430}  & \inred{5} & \inred{221} & \yes & \inred{431}  & \inred{3} \\
    L2 Distance 32  & \inred{262} & \no  & \inred{790}  & \inred{6} & \inred{229} & \yes & \inred{772}  & \inred{4} \\
    Hamm. Dist. 16  & \inred{247} & \no  & \inred{980}  & \inred{5} & \inred{226} & \yes & \inred{951}  & \inred{4} \\
    Hamm. Dist. 32  & \inred{271} & \no  & \inred{1490} & \inred{7} & \inred{228} & \yes & \inred{1402} & \inred{4} \\
    Poly. Reg. 8    & \inred{219} & \yes & \inred{315}  & \inred{2} & \inred{212} & \yes & \inred{309}  & \inred{2} \\
    Linear Reg. 8   & \inred{188} & \yes & \inred{140}  & \inred{2} & \inred{188} & \yes & \inred{140}  & \inred{1} \\
    Box Blur 4      & \inred{205} & \yes & \inred{96}   & \inred{3} & \inred{203} & \yes & \inred{92}   & \inred{2} \\
    Mat. Mul. 4     & \inred{252} & \no  & \inred{512}  & \inred{4} & \inred{226} & \yes & \inred{498}  & \inred{3} \\
    Gx 4            & \inred{214} & \yes & \inred{530}  & \inred{3} & \inred{210} & \yes & \inred{521}  & \inred{2} \\
    Max 4           & \inred{233} & \no  & \inred{1210} & \inred{5} & \inred{225} & \yes & \inred{1188} & \inred{4} \\
    Sort 3          & \inred{239} & \no  & \inred{1380} & \inred{5} & \inred{229} & \yes & \inred{1366} & \inred{4} \\
    tree-100-50-5   & \inred{266} & \no  & \inred{2010} & \inred{6} & \inred{228} & \yes & \inred{1985} & \inred{4} \\
    \midrule
    Violations      & \multicolumn{4}{c}{\inred{10/14}} & \multicolumn{4}{c}{\inred{0/14}} \\
    \bottomrule
  \end{tabular}
\end{table}

Figure~\ref{fig:exp-unified-vs-original} and Table~\ref{tab:exp-unified-vs-original} show the results. CHEHAB RL exceeds the budget on \inred{10} of the \inred{14} benchmarks, by up to \inred{17.8\%} (\inred{Hamm. Dist. 32}). \ToolName{} satisfies the budget on all benchmarks. The constraint does not cost execution time: \ToolName{} is as fast or faster on \inred{13} benchmarks, with a geometric-mean speedup of \inred{$1.02\times$}. It also needs fewer keys: \inred{44} keys (\inred{792}~MB) in total, against \inred{64} (\inred{1152}~MB) for CHEHAB RL, a reduction of \inred{31\%}. We attribute this difference to when keys are handled. CHEHAB RL fixes the vectorization first and then replaces rotations by chains of rotations, which adds operations and noise. \ToolName{} accounts for keys and noise while it chooses the vectorization.





\subsubsection{Joint Optimization vs. Separate Optimization
\label{sec:evaluation:joint-vs-separate}}



\subsubsection{Trade-off Evolution Across Preference Weights
\label{sec:evaluation:tradeoff-evolution}}



\subsubsection{Multi-Objective Optimization Quality Across Benchmarks
\label{sec:evaluation:mo-quality}}




\subsubsection{Constraint Satisfaction Under Different Noise Budgets
\label{sec:evaluation:constraint-satisfaction}}



\subsubsection{Noise Estimator Accuracy
\label{sec:evaluation:noise-estimator}}



\subsubsection{Stress Test on Threshold Benchmarks
\label{sec:evaluation:stress-test}}





\subsection{Ablation Studies
\label{sec:evaluation:ablations}}

\subsubsection{Impact of Bisection on Pareto Front Quality
\label{sec:evaluation:ablation-bisection}}




\subsubsection{Impact of Biased Preference Sampling
\label{sec:evaluation:ablation-biased-sampling}}



\subsubsection{Impact of Safety Rollback
\label{sec:evaluation:ablation-safety-rollback}}




\subsubsection{Impact of Constraint Enforcement Mechanisms
\label{sec:evaluation:ablation-constraint-mechanisms}}




% \shrink[1]
% \subsection{Additional Evaluations
% \label{sec:evaluation:additional-evaluations}}

% \placeholder[14]
